\scanpage{58}

\diagram{\begin{tikzpicture}
\node[anchor=north west,inner sep=0pt] (g) at (0,0) {\begin{tikzcd}[row sep=large]
\alpha_i\text{ bijectif pour }m\text{ grand}\arrow[d,Leftrightarrow]\\
\operatorname{Supp}E^{r-i}\subset Z\arrow[d,Rightarrow]\\
H^i_x(F_x)=0\quad\text{pour }\forall\ x\text{ fermé \uline{en $U$}.}
\end{tikzcd}};
\draw (g.north west) .. controls +(-.3,0) and +(-.2,.3) .. (g.west)
 .. controls +(-.2,-.3) and +(-.3,0) .. (g.south west);
\end{tikzpicture}}

\diagram{\begin{tikzpicture}
\node[anchor=north west,inner sep=0pt] (g) at (0,0) {\begin{tikzcd}[row sep=large]
\alpha_i\text{ bijectif pour }m\text{ grand},\quad i\leq n\arrow[d,Leftrightarrow]\\
\operatorname{supp}E^{r-i}\subset Z\quad\text{pour }i\leq n\arrow[d,Leftrightarrow]\\
\operatorname{prof}(F_x)>n\quad\text{pour }\forall\ x\text{ fermé \uline{en} }U
\arrow[d,Leftrightarrow]\\
\begin{gathered}
H^i(U,F(-m))=\deleted{H^{\unsure{0|i}}(\unsure{X|U},F(-m))}
\inserted{0\quad\text{pour }m\text{ grand},\ i<n\quad\text{\uuline{et}}}
\quad\deleted{\text{bijectif pour }i\leq n}\\
H^n(U,F(-m))\longrightarrow H^{n+1}_{\inserted{Z}}(X,F(-m))
\quad\deleted{\text{surjectif pour }i\leq n}\ \inserted{\text{injectif}}
\end{gathered}
\end{tikzcd}};
\draw (g.north west) .. controls +(-.3,0) and +(-.2,.35) .. (g.west)
 .. controls +(-.2,-.35) and +(-.3,0) .. (g.south west);
\end{tikzpicture}}

\diagram{\begin{tikzpicture}
\node[anchor=north west,inner sep=0pt] (g) at (0,0) {\begin{tikzcd}[row sep=large]
\begin{gathered}
\alpha_i\text{ bijectif pour }m\text{ grand},\quad i\leq n\\
\text{et }\alpha_i\text{ injectif pour }m\text{ grand},\quad i=n+1
\end{gathered}\arrow[d,Leftrightarrow]\\
\begin{gathered}
\deleted{\unsure{\operatorname{prof}}}\ \operatorname{supp}E^{r-i}\subset Z
\quad\text{pour }i\leq n,\quad\operatorname{supp}E^{r-n-1}\ \text{\unsure{ne peut}}\\
\text{\unsure{avoir un pt de $Z$ pour}}\ \text{\illegible{2}}\ \text{\unsure{pts isolés}}
\end{gathered}\arrow[d,Leftrightarrow]\\
\begin{gathered}
\operatorname{prof}(F_x)>n\text{ si }x\text{ fermé en }U,\quad\text{et}\\
\operatorname{prof}(F_x)>n+1\text{ si de plus }x\in V\cap U,\quad V\text{ voisinage}\\
\text{ouvert convenable de }Z
\end{gathered}\arrow[d,Leftrightarrow]\\
H^i(U,F(-m))=0\quad\text{pour }m\text{ grand},\quad i\leq n\arrow[d,Leftrightarrow]\\
\coprod_{m\unsure{\geq|\gg}0}H^i(U,F(m))\text{ de type fini sur }S(\unsure{-h})\quad i\leq n
\end{tikzcd}};
\draw (g.north west) .. controls +(-.4,0) and +(-.25,.5) .. (g.west)
 .. controls +(-.25,-.5) and +(-.4,0) .. (g.south west);
\node[draw,circle] at (-1.2,-3) {$B_n$};
\node[draw,ellipse,align=center] (m) at (-2,-7.1)
 {$H^i(U,F)$\\de dim finie\\pour $i\leq n$};
\draw[->,double] (-.2,-6.65) -- node[above] {cf $2^\circ$} (m.north east);
\end{tikzpicture}}

Cela \unsure{résout} la question (a). Sous les conditions de (*), on a donc
[si $t$ est $F$-régulier]
\[
\left\{\begin{gathered}
H^i(U,F)\overset{\sim}{\longrightarrow}H^i(\hat X,\hat F)
\quad\text{est}\quad
\begin{cases}\text{iso si }i<n\\\text{mono si }i=n\end{cases}\\
H^i(\hat X,\hat F)\simeq\varprojlim H^i(X_m,F_m)\quad\text{pour }i\leq n
\end{gathered}\right.
\]

\margin{57}
